\documentclass[10pt,a4paper]{amsart}
\usepackage[margin=19mm]{geometry}
\usepackage{amsmath,amssymb,mathtools}
\usepackage[scr=rsfs,cal=euler]{mathalfa}
\usepackage{xfrac}
\usepackage[unicode,hidelinks,bookmarks=false]{hyperref}
\newcommand{\ie}{i.e.\ }
\newcommand{\cf}{cf.\ }
\newcommand{\resp}{respectively\ }
\newcommand{\Subsection}{\subsection}
\providecommand{\nobreakdash}{\nobreak\hbox{-}}
\setlength{\emergencystretch}{2em}
\multlinegap0pt
\usepackage{amssymb}
\usepackage{mathtools}
\usepackage[shortlabels]{enumitem}
\newtheorem{proposition}{Proposition}
\title{Holographic functions and neural networks}
\author{Balazs Szegedy}
\date{Source: arXiv:2605.22666v1 (CC BY 4.0)}
\begin{document}
\maketitle

\noindent\emph{Source and licence.} Holographic functions and neural networks. Version arXiv:2605.22666v1 (CC BY 4.0), distributed by arXiv under the Creative Commons Attribution 4.0 International licence, http://creativecommons.org/licenses/by/4.0/, as recorded in the arXiv licence metadata for this exact version and shown on the abstract page https://arxiv.org/abs/2605.22666v1. Full licence notice: \texttt{licenses/arxiv-2605.22666-CC-BY-4.0.txt}.\par\smallskip

\noindent\textit{Editorial scope.} Counted unit: the article's proposition prop:identical-nonidentical-equivalence (source0.tex lines 772--795). The article's complete original proof of it is delivered with it (source0.tex lines 797--1009, one \texttt{proof} environment of 213 lines). No mathematics is written, repaired or re-derived: the statement and the proof are the article's own lines, in the article's own notation, and no retained line is substituted or re-flowed.  No further source-local context is needed: every cross-reference of every retained line resolves inside the two delivered spans.  Nothing was omitted from any retained span, so the package carries no omission record.  No retained line contains a citation, so the wrapper carries no bibliography and no bibliography entry.  No retained line contains a figure environment, an image-inclusion command or drawing code, so no image is delivered and the package declares no asset dependency.  Editorial formatting only, in addition: an AMS \texttt{amsart} preamble that restates the article's own declarations and macros used by the retained lines, namely the environments \texttt{proposition} on separate counters, together with the standard packages those lines need.  The retained environments print in this document as Proposition 1 in order of appearance, whereas the article prints the same items with the numbering of its own full document.  The complete native source of the article as distributed in the arXiv e-print for this exact version is bundled unchanged as \texttt{sources/201148/source0.tex}.
\begin{proposition}\label{prop:identical-nonidentical-equivalence}
Let \(k\in\mathbb N\) and \(0<\varepsilon<1\).  Put
\[
        \alpha=\frac{\varepsilon^2}{4},
        \qquad
        \eta=\frac{\varepsilon^2}{2}.
\]
Suppose that \(f:\{0,1\}^n\to[0,1]\) is
\((k,\alpha)\)-holographic with non-identical sampling.  Then \(f\) is
identically sampled \((r,\varepsilon)\)-holographic for every integer
\(r\) satisfying
\[
        k e^{-r/k}\leq \eta .
\]
In particular, it is enough to take
\[
        r\geq k\log\frac{k}{\eta}
        =
        k\log\frac{2k}{\varepsilon^2}.
\]
Consequently, the identically sampled and non-identically sampled versions of
holography define the same qualitative notion after changing the complexity
and accuracy parameters.
\end{proposition}

\begin{proof}
Assume that \(f\) is \((k,\alpha)\)-holographic with non-identical sampling.
Thus there are probability measures
\[
        \mu_1,\ldots,\mu_k
\]
on \([n]\), and test functions
\[
        f_s:\{0,1\}^k\to[0,1],
        \qquad s\in[n]^k,
\]
such that for every \(x\in\{0,1\}^n\),
\[
        \mathbb P\left(
        \left|f(x)-f_S(x_{S_1},\ldots,x_{S_k})\right|
        \leq \alpha
        \right)
        \geq 1-\alpha ,
\]
where \(S_i\) are independent and \(S_i\sim\mu_i\).

Define the average measure
\[
        \mu=\frac1k\sum_{i=1}^k \mu_i .
\]
A sample from \(\mu\) may be generated in two steps: first choose a hidden label
\(I\in[k]\) uniformly, and then sample from \(\mu_I\).  More precisely, let
\[
        I_1,\ldots,I_r
\]
be independent uniform random variables on \([k]\), and, conditional on these
labels, let
\[
        T_j\sim \mu_{I_j}
\]
independently for \(j=1,\ldots,r\).  Then the variables \(T_1,\ldots,T_r\) are
independent and each has law \(\mu\).

Let \(E\) be the event that every label \(1,\ldots,k\) appears among
\(I_1,\ldots,I_r\).  By the union bound,
\[
        \mathbb P(E^c)
        \leq
        k\left(1-\frac1k\right)^r
        \leq
        k e^{-r/k}
        \leq \eta .
\]

On the event \(E\), define \(\tau_i\) to be the first index \(j\) such that
\(I_j=i\).  Thus
\[
        I_{\tau_i}=i,
        \qquad i=1,\ldots,k.
\]
The selected coordinates
\[
        T_{\tau_1},\ldots,T_{\tau_k}
\]
are independent and have respective laws
\[
        \mu_1,\ldots,\mu_k.
\]
Indeed, conditional on any label sequence for which \(E\) holds, the selected
coordinate \(T_{\tau_i}\) is sampled from \(\mu_i\), and the selected
coordinates for distinct labels are independent.

For a fixed input \(x\), define the labelled reconstruction
\[
        Y
        =
        f_{(T_{\tau_1},\ldots,T_{\tau_k})}
        (x_{T_{\tau_1}},\ldots,x_{T_{\tau_k}})
\]
on the event \(E\).  On \(E^c\), define \(Y=0\), say.  Since on \(E\) the
selected coordinates have the same joint distribution as the original
non-identically sampled tuple, we have
\[
        \mathbb P\left(|Y-f(x)|>\alpha\right)
        \leq
        \alpha+\eta .
\]
Because \(Y\in[0,1]\) and \(f(x)\in[0,1]\), it follows that
\[
        \mathbb E |Y-f(x)|
        \leq
        \alpha+\mathbb P(|Y-f(x)|>\alpha)
        \leq
        2\alpha+\eta .
\]

So far the reconstruction used the hidden labels \(I_1,\ldots,I_r\).  The
definition of identical sampling, however, only permits the test function to
depend on the observed locations \(T_1,\ldots,T_r\) and on the observed bits.
We now remove the labels by conditional expectation.  This posterior averaging
is only part of the construction of the deterministic test functions; the
eventual identical-sampling procedure itself samples only the locations from
the single measure \(\mu\).

For \(t=(t_1,\ldots,t_r)\in[n]^r\), define a probability distribution on label
tuples \(i=(i_1,\ldots,i_r)\in[k]^r\) by the posterior rule
\[
        \mathbb P(I_1=i_1,\ldots,I_r=i_r\mid T_1=t_1,\ldots,T_r=t_r)
        =
        \prod_{j=1}^r
        \frac{\mu_{i_j}(t_j)}
        {\sum_{\ell=1}^k \mu_\ell(t_j)}
\]
whenever the denominators are nonzero.  This is just Bayes' formula with
respect to the auxiliary hidden-label construction: the factor
\((1/k)\mu_{i_j}(t_j)\) in the joint law is divided by the mixture mass
\(\mu(t_j)=(1/k)\sum_{\ell=1}^k\mu_\ell(t_j)\), so the factors \(1/k\) cancel.
If a denominator is zero, the value is irrelevant because such a location is
never sampled from \(\mu\), and we may choose the conditional distribution
arbitrarily.

Given \(t\in[n]^r\) and \(a=(a_1,\ldots,a_r)\in\{0,1\}^r\), define
\[
        g_t(a)
\]
to be the conditional expectation, over these posterior labels, of the labelled
reconstruction described above, with the observed bits \(a_j\) substituted for
\(x_{t_j}\).  More explicitly, for each label tuple \(i\), if every label
\(1,\ldots,k\) occurs among \(i_1,\ldots,i_r\), let \(\tau_\ell(i)\) be the
first index \(j\) with \(i_j=\ell\), and set
\[
        Y_i(t,a)
        =
        f_{(t_{\tau_1(i)},\ldots,t_{\tau_k(i)})}
        (a_{\tau_1(i)},\ldots,a_{\tau_k(i)}).
\]
If not every label occurs, set \(Y_i(t,a)=0\).  Then define
\[
        g_t(a)
        =
        \mathbb E\bigl[ Y_I(t,a)\mid T=t\bigr].
\]
This is a deterministic function
\[
        g_t:\{0,1\}^r\to[0,1].
\]

For the actual input \(x\), the random variable
\[
        g_T(x_{T_1},\ldots,x_{T_r})
\]
is exactly
\[
        \mathbb E(Y\mid T_1,\ldots,T_r,x_{T_1},\ldots,x_{T_r}).
\]
Therefore, by Jensen's inequality,
\[
\begin{aligned}
        \left|
        g_T(x_{T_1},\ldots,x_{T_r})-f(x)
        \right|
        &=
        \left|
        \mathbb E(Y-f(x)\mid T_1,\ldots,T_r,x_{T_1},\ldots,x_{T_r})
        \right|                                                   \\
        &\leq
        \mathbb E\left(
        |Y-f(x)|
        \mid T_1,\ldots,T_r,x_{T_1},\ldots,x_{T_r}
        \right).
\end{aligned}
\]
Taking expectations gives
\[
        \mathbb E
        \left|
        g_T(x_{T_1},\ldots,x_{T_r})-f(x)
        \right|
        \leq
        \mathbb E |Y-f(x)|
        \leq
        2\alpha+\eta .
\]
By Markov's inequality,
\[
        \mathbb P\left(
        \left|
        g_T(x_{T_1},\ldots,x_{T_r})-f(x)
        \right|
        >\varepsilon
        \right)
        \leq
        \frac{2\alpha+\eta}{\varepsilon}.
\]
With our choices
\[
        \alpha=\frac{\varepsilon^2}{4},
        \qquad
        \eta=\frac{\varepsilon^2}{2},
\]
we have
\[
        2\alpha+\eta=\varepsilon^2.
\]
Hence
\[
        \mathbb P\left(
        \left|
        g_T(x_{T_1},\ldots,x_{T_r})-f(x)
        \right|
        >\varepsilon
        \right)
        \leq \varepsilon .
\]
This holds for every \(x\in\{0,1\}^n\).  Therefore \(f\) is
identically sampled \((r,\varepsilon)\)-holographic with identical sampling
from the single measure \(\mu\).
\end{proof}

\end{document}
